% 原书第 214—215 页；第 13.7 节与致谢全文。
\section{讨论}
\label{sec:ch13-discussion}

本章建立在随机试验中协变量调整的丰富研究历史之上，例如 Fisher（1932）、Cochran（1957）、Cox and McCullagh（1982）、Tsiatis et al.（2008）、Zhang et al.（2008）、Moore et al.（2011）、Yuan et al.（2012）、Shen et al.（2014）以及 Colantuoni and Rosenblum（2015）。具体而言，Rubin and van der Laan（2008）提出经验效率最大化原则，作为选择条件结局均值 $\bar Q_0(A,W)$ 估计量的一种策略，使估计有效影响曲线的经验方差最小。不过，他们的方法需要求解一个加权非线性最小二乘问题；我们的方法只要求研究者计算样本方差。

C-TMLE 的近期发展提出，通过协同估计暴露机制，在观察性研究的 TMLE 靶向步骤中，尽可能多地减小偏差。在随机试验中，回归模型的设定错误不会带来偏差风险（Rosenblum and van der Laan，2010b）。因此，本章实施的协同方法，通过估计已知的暴露机制，只用于提高精度。本章将这一方案推广至有配对匹配和无配对匹配的随机试验，用于总体平均处理效应及样本平均处理效应的估计与推断。因此，这一方法消除了人们普遍担忧的配对匹配试验“分析上的局限”，例如 Klar and Donner（1997）、Imbens（2011）以及 Campbell（2014）。由于逐步实施的算法，包括算法库的定义，已经预先指定，就不会因临时作出的分析决定，而产生偏差或误导性的推断。此外，将未经调整的估计量纳入候选库，就不再需要关于是否应调整协变量的指南，例如 Moore et al.（2011）以及 Colantuoni and Rosenblum（2015）。最后，这一方法针对科学问题，即总体或样本效应，以及研究设计，即是否采用配对匹配来制定。关于是否调整、如何调整的决定，都通过严格、有原则的方法作出，从而减少了统计学中的一些“人为技巧”。

\noindent\textbf{致谢。}
本章报告的研究得到美国国立卫生研究院（NIH）下属美国国家过敏与传染病研究所（NIAID）艾滋病部的资助，项目编号为 R01-AI074345、R37-AI051164、UM1AI069502 和 U01AI099959。本章内容由作者自行负责，并不必然代表 NIH 的官方观点。

